% GENERATED FILE. Edit canonical lessons, then run npm run generate:books.
% canonical-source-hash: fnv1a64:ad1ff4e2c888329f
% canonical-lessons: GU-C104-anubhav, GU-C104-tev, GU-C104-adat, GU-C104-layakat, GU-C104-kala

\chapter{Experience, Habit, Habit, Qualification, Art}
\label{ch:gu-experience-habit-habit-qualification-art}
\hlchaptermodality{104}

\begin{chapteropening}
\textbf{By the end of this chapter:} \emph{I can say experience, a habit, a habit, a qualification and an art.}

Complete the last lesson of chapter 104: I can say experience, a habit, a habit, a qualification and an art.
\end{chapteropening}

\section[\texorpdfstring{\gu{અનુભવ} (anubhav)}{અનુભવ (anubhav)}]{\gu{અનુભવ} (anubhav) — experience}

\label{lesson:GU-C104-anubhav}



\begin{quote}
Before the new one: say the Gujarati for wealth, then the Gujarati for an owner.
\end{quote}

\subsection*{You'll want to know: \gu{અનુભવ}}
\textbf{\gu{અનુભવ}} (\emph{anubhav}) — \textquotedblleft{}experience\textquotedblright{}.

\begin{cousinweb}
\textbf{\gu{અનુભવ}} (\emph{anubhav}) is Sanskrit \emph{anubhava}, \textquotedblleft{}perception, experience\textquotedblright{}, taken over whole as a learned loan, a tatsama.
\end{cousinweb}

\begin{grammarlens}[title={What you've built}]
One more word for this chapter.
\end{grammarlens}

\subsection*{Your turn}
\begin{itemize}
\raggedright
  \item \emph{Say it:} \emph{anubhav}
  \item \emph{Say it:} \emph{anubhav}, once more
  \item \emph{Say it:} say \emph{anubhav}
  \item \emph{Recall:} say the Gujarati for wealth, then the Gujarati for an owner, then say \emph{anubhav} again
\end{itemize}

\begin{tcolorbox}[breakable,colback=teal!4,colframe=teal!35!black,title={Before you move on}]
What does \gu{અનુભવ} mean? (\textquotedblleft{}Experience\textquotedblright{}.) Say it once more.
\end{tcolorbox}
\section[\texorpdfstring{\gu{ટેવ} (ṭev)}{ટેવ (ṭev)}]{\gu{ટેવ} (ṭev) — a habit}

\label{lesson:GU-C104-tev}



\begin{quote}
Before the new one: say the Gujarati for an owner, then the Gujarati for experience.
\end{quote}

\subsection*{You'll want to know: \gu{ટેવ}}
\textbf{\gu{ટેવ}} (\emph{ṭev}) — \textquotedblleft{}a habit\textquotedblright{}.

\begin{cousinweb}
\textbf{\gu{ટેવ}} (\emph{ṭev}) is a tadbhava, a word inherited and reshaped by sound change. It comes from \emph{*ṭev}, \textquotedblleft{}appoint, fix\textquotedblright{}, a form CDIAL reconstructs from its descendants (CDIAL 5475). CDIAL lists Gujarati \emph{ṭev}, \textquotedblleft{}habit\textquotedblright{}, under it.
\end{cousinweb}

\begin{grammarlens}[title={What you've built}]
One more word for this chapter.
\end{grammarlens}

\subsection*{Your turn}
\begin{itemize}
\raggedright
  \item \emph{Say it:} \emph{ṭev}
  \item \emph{Say it:} \emph{ṭev}, once more
  \item \emph{Say it:} say \emph{ṭev}
  \item \emph{Recall:} say the Gujarati for an owner, then the Gujarati for experience, then say \emph{ṭev} again
\end{itemize}

\begin{tcolorbox}[breakable,colback=teal!4,colframe=teal!35!black,title={Before you move on}]
What does \gu{ટેવ} mean? (\textquotedblleft{}A habit\textquotedblright{}.) Say it once more.
\end{tcolorbox}
\section[\texorpdfstring{\gu{આદત} (ādat)}{આદત (ādat)}]{\gu{આદત} (ādat) — a habit, a custom}

\label{lesson:GU-C104-adat}



\begin{quote}
Before the new one: say the Gujarati for experience, then the Gujarati for a habit.
\end{quote}

\subsection*{You'll want to know: \gu{આદત}}
\textbf{\gu{આદત}} (\emph{ādat}) — \textquotedblleft{}a habit, a custom\textquotedblright{}.

\begin{cousinweb}
\textbf{\gu{આદત}} (\emph{ādat}) is a loan from Arabic \emph{ʿāda}, \textquotedblleft{}custom, habit\textquotedblright{}, which reached Gujarati through Persian. Its Arabic root is \emph{ʿ-w-d} (Lane).
\end{cousinweb}

\begin{grammarlens}[title={What you've built}]
One more word for this chapter.
\end{grammarlens}

\subsection*{Your turn}
\begin{itemize}
\raggedright
  \item \emph{Say it:} \emph{ādat}
  \item \emph{Say it:} \emph{ādat}, once more
  \item \emph{Say it:} say \emph{ādat}
  \item \emph{Recall:} say the Gujarati for experience, then the Gujarati for a habit, then say \emph{ādat} again
\end{itemize}

\begin{tcolorbox}[breakable,colback=teal!4,colframe=teal!35!black,title={Before you move on}]
What does \gu{આદત} mean? (\textquotedblleft{}A habit, a custom\textquotedblright{}.) Say it once more.
\end{tcolorbox}
\section[\texorpdfstring{\gu{લાયકાત} (lāyakāt)}{લાયકાત (lāyakāt)}]{\gu{લાયકાત} (lāyakāt) — a qualification}

\label{lesson:GU-C104-layakat}



\begin{quote}
Before the new one: say the Gujarati for a habit, then the Gujarati for a habit.
\end{quote}

\subsection*{You'll want to know: \gu{લાયકાત}}
\textbf{\gu{લાયકાત}} (\emph{lāyakāt}) — \textquotedblleft{}a qualification\textquotedblright{}.

\begin{cousinweb}
\textbf{\gu{લાયકાત}} (\emph{lāyakāt}) is a loan, usually traced to Persian \emph{liyāqat}, \textquotedblleft{}fitness, worth\textquotedblright{}, from Arabic, beside \emph{lāyak}, \textquotedblleft{}fit\textquotedblright{}. The Arabic root is \emph{l-y-q} (Lane).
\end{cousinweb}

\begin{grammarlens}[title={What you've built}]
One more word for this chapter.
\end{grammarlens}

\subsection*{Your turn}
\begin{itemize}
\raggedright
  \item \emph{Say it:} \emph{lāyakāt}
  \item \emph{Say it:} \emph{lāyakāt}, once more
  \item \emph{Say it:} say \emph{lāyakāt}
  \item \emph{Recall:} say the Gujarati for a habit, then the Gujarati for a habit, then say \emph{lāyakāt} again
\end{itemize}

\begin{tcolorbox}[breakable,colback=teal!4,colframe=teal!35!black,title={Before you move on}]
What does \gu{લાયકાત} mean? (\textquotedblleft{}A qualification\textquotedblright{}.) Say it once more.
\end{tcolorbox}
\section[\texorpdfstring{\gu{કળા} (kaḷā)}{કળા (kaḷā)}]{\gu{કળા} (kaḷā) — an art}

\label{lesson:GU-C104-kala}



\begin{quote}
Before the new one: say the Gujarati for a habit, then the Gujarati for a qualification.
\end{quote}

\subsection*{You'll want to know: \gu{કળા}}
\textbf{\gu{કળા}} (\emph{kaḷā}) — \textquotedblleft{}an art\textquotedblright{}.

\begin{cousinweb}
\textbf{\gu{કળા}} (\emph{kaḷā}) is Sanskrit \emph{kalā}, \textquotedblleft{}any practical art\textquotedblright{}, lightly reshaped with Gujarati \emph{ḷ}, a semi-tatsama. CDIAL takes the Sanskrit word from Dravidian. The only Gujarati word it lists there is \emph{kaḷ}, \textquotedblleft{}door-handle, gadget, scheme\textquotedblright{} (CDIAL 2927).
\end{cousinweb}

\begin{grammarlens}[title={What you've built}]
One more word for this chapter.
\end{grammarlens}

\subsection*{Your turn}
\begin{itemize}
\raggedright
  \item \emph{Say it:} \emph{kaḷā}
  \item \emph{Say it:} \emph{kaḷā}, once more
  \item \emph{Say it:} say \emph{kaḷā}
  \item \emph{Recall:} say the Gujarati for a habit, then the Gujarati for a qualification, then say \emph{kaḷā} again
\end{itemize}

\begin{tcolorbox}[breakable,colback=teal!4,colframe=teal!35!black,title={Before you move on}]
What does \gu{કળા} mean? (\textquotedblleft{}An art\textquotedblright{}.) Say it once more.
\end{tcolorbox}
